\scanpage{23}\mspage{21}
Il est immédiat que si $E$ est admissible, $E'$ aussi (et
inversement). Ceci dit, on a le

\uline{Lemme 2} \quad $E$ et $E'$ étant comme ci-dessus, il existe un
\unsure{hom|Hom} \deleted{\unsure{fonct.}} de foncteurs entre $\varphi_E$ et $\varphi_{E'}$ (je ne dis
pas dans quel sens va la flèche !) \deleted{\unsure{$\leftarrow$}}.

\uline{Cor} \quad Énoncé dual, quand $E'$ est déduit de $E$ en translatant
à droite \inserted{d'un cran} un intervalle \inserted{connexe} de longueur paire \inserted{$I=[r,s]$} de $E$, \deleted{tel que l'on ait}
$s\ne 2k$ et $s+1\notin E$.

En effet, dans ce cas $E$ est déduit de $E'$ par une \unsure{op.} du
type explicité dans le lemme 2.

La connexité de la catégorie $\underline{\mathrm{Hom}}^{!}(\unsure{\mathscr{L}|\mathcal{L}}_k,\mathscr{L}_l)$ est
maintenant claire. Si $k$ et $l$ ont même parité,
en poussant les comp. connexes de $E$ vers la gauche,
on met $E$ sous la forme ``normale''
\[
E_0=[1,2(k-l)]\quad .
\]
Si $k$ et $l$ sont de parité inverse, on \deleted{\unsure{ne}} ne peut
pousser \deleted{tout} la comp. conn. de $E$ qui contient
l'élément dernier $2k$, car elle est de cardinal impair.
Il faut laisser $2k$ dans $E$ et pousser le reste, et
on ramène $E$ à la forme normale
\[
E_0=[1,2(k-l)-1]\cup\{2k\}\quad ,
\]
cqfd !

\deleted{\unsure{§}} Ce qui précède justifie notre raisonnement,
qui donnait par composition des chemins des
applications de composition bien définies
\[
\underbrace{\operatorname{Hom}_{\unsure{\Sigma|\overline{Z}|\overline{\Sigma}|Z}}(x,y)}_{\substack{\Vert\\\pi_0\underline{\mathrm{Ch}}_\Sigma(x,y)}}
\times\underbrace{\operatorname{Hom}_{\Sigma}(y,z)}_{\substack{\Vert\\\pi_0\underline{\mathrm{Ch}}_\Sigma(y,z)}}
\unsure{\to|\mathord{-}}
\underbrace{\operatorname{Hom}_{\Sigma}(x,z)}_{\substack{\Vert\\\pi_0\underline{\mathrm{Ch}}_\Sigma(x,z)}}
\]
% archive sheet number: 22
